%% kapa: exam=1 round=28 year=2017 period=1 form=A subject=civil_law subject_ko=민법 seq=30 number=30
%% source: Q-Net 감정평가사 기출문제 (cst00309) artlSeq=5199603 file=extracted/1008/5199603/2017년 제28회 감정평가사 1차시험 문제지/1교시 A형.pdf
%% conversion: kordoc 4.24.51 -> tools/gen1.py (auto); review=proofread-sonnet; flags=none
\begin{question}{30}
첨부에 관한 설명으로 옳지 않은 것은? (다툼이 있으면 판례에 따름)
\begin{choices}
  \item 주종의 구별이 있는 동산과 동산이 부합된 합성물은 주된 동산의 소유자에게 속한다.
  \item 완성된 건물은 토지에 부합하지 않는다.
  \item 가공물은 원칙적으로 원재료 소유자에게 속한다.
  \item 부동산에 부합되어 동산의 소유권이 소멸하는 경우, 그 동산을 목적으로 한 질권은 소멸하지 않는다.
  \item 첨부에 의해 손해를 받은 자는 부당이득에 관한 규정에 의하여 보상을 청구할 수 있다.
\end{choices}
\end{question}
\answer{4}
